\section{Introduction}

Gallbladder diseases---cholelithiasis, cholecystitis, polyps, adenomyomatosis, carcinoma, GBC---affect hundreds of millions worldwide. GBC is especially deadly when caught late; five-year survival drops below 20\% \cite{ref1}, \cite{ref2}. Ultrasound is the go-to first-line scan for the gallbladder. Yet a systematic review and meta-analysis shows its sensitivity for acute cholecystitis sits at a mere 81\% (studies range from 50\% to 100\%) and specificity at 83\% (33\% to 100\%). That spread mirrors real-world operator and protocol variation \cite{ref3}. Such inconsistency pushes us toward computer-aided classification, a decision-support tool, while also highlighting the promise and the caution that accompany AI in medicine \cite{ref4}, \cite{ref5}, \cite{ref6}. CNNs now match or outshine earlier gallbladder-ultrasound classifiers \cite{ref7}, but clinical uptake stalls at opacity. A CNN is a black-box: it spits out a verdict without showing why. Radiologists need visual proof that the model is looking at real anatomy---wall irregularity, acoustic shadowing, lumen echogenicity---rather than noise or artefacts \cite{ref9}, \cite{ref10}.

XAI techniques---Grad-CAM \cite{ref11}, Grad-CAM++ \cite{ref12}, gradient saliency maps---offer class-discriminative heatmaps. A heatmap that appears anatomically plausible is not proof of reliability. Ghassemi et~al. \cite{ref9} and Rudin \cite{ref8} note that most clinical XAI studies stop at visual inspection, never checking if the explanation actually points to the right area, survives small changes, reflects the true driver of the prediction, or repeats across separate training runs. Sanity-check studies \cite{ref13} and interpretability benchmarks \cite{ref14} reveal that some popular saliency methods fail these tests yet still look convincing. A convincing heatmap alone is insufficient; we must quantify reliability.

Earlier gallbladder-ultrasound studies have employed XAI, and at least one provided clinician-grounded, quantitative localization evidence \cite{ref19}. Yet a structured, non-systematic literature review (Section~\ref{sec:related-work}) found no peer-reviewed work that pulls together a multi-method post-hoc attribution comparison with quantitative metrics for cross-training consistency, perturbation-based faithfulness, explanation stability. And sensitivity to learned parameters---all within one framework. When several XAI methods appear together in this field, they tend to be compared qualitatively, not measured along these reliability dimensions. This paper asks: how much does the choice of XAI method influence a nine-class gallbladder-ultrasound classifier. And are the explanations reproducible or just artifacts of a single training run? The goal is not to introduce XAI to gallbladder ultrasound---previous work has done that---but to quantitatively test whether common post-hoc visual explanations are reliable. We tackle this with a multi-method, quantitatively evaluated XAI analysis, confined to a single public dataset and evaluated at the image level. The dataset's original publication says the images came from three to four hospitals in Baghdad, Iraq. But the public release strips away hospital, patient, or scanner IDs, so we cannot recover or stratify site provenance (Section~\ref{sec:methods}.1). We make the scope clear: this is an internal-validation study on data with unknown provenance, not a claim of external multicentre generalizability or clinical deployment. Those larger questions will guide the validation programme we start here (Section~\ref{sec:future-research}).

Our work has five aims. First, we train and test a ResNet-50 to perform nine-class gallbladder-ultrasound classification. Second, we apply five XAI methods---Grad-CAM, Grad-CAM++, Saliency Maps, Eigen-CAM, Score-CAM---to the same trained model. So we can compare them head-to-head instead of looking at one in isolation. Third, we calculate a cosine-based consistency score to see if each method's explanation stays the same across independently trained model replicas. Fourth, we test faithfulness by deletion/insertion against a random control. Fifth, we assess stability under clinically invariant image transforms and check if the explanations hinge on learned weights rather than just the architecture, using a parameter-randomisation sanity check. Alongside, we calibrate the classifier's predicted probabilities with temperature scaling, so the XAI outputs can be paired with trustworthy confidence values.

This paper offers six contributions. First, a two-stage leakage audit of the UIdataGB release itself: an initial split de-duplicated only by single-orientation perceptual hashing looked clean but still let near-identical frames from the same imaging session, and rotation-augmented duplicates the dataset's own release pipeline introduces, cross the train/validation boundary; correcting both drops mean accuracy from an inflated 93.05\% to an honestly-measured 72.56\%~$\pm$~2.35\% across four independently seeded models on a genuine, cluster-aware held-out test set. We treat this correction, and the willingness to report the lower number, as a contribution in its own right, not a footnote. Second, a quantitative XAI benchmark that runs five methods---Grad-CAM, Grad-CAM++, Saliency Maps, Eigen-CAM, Score-CAM---on a ResNet-50 gallbladder-ultrasound model, covering more methods and classes than any previous study (see Table~\ref{tbl:9}). Third, a quantitative XAI-reliability suite: cosine-based cross-replicate consistency $O^c$, faithfulness tests, stability checks, a model-randomisation sanity check, calibration, and a shortcut audit, together with a null-control test of the consistency metric itself against untrained-model and synthetic-heatmap floors, which the metric had not previously been checked against. This moves us past ``heatmaps that look right'' to explanations that are actually measured to be reliable. Fourth, a failure-case analysis grounded in clinician-affirmed evidence that looks at specific high-confidence misclassifications instead of just overall accuracy. Fifth, an honest validation roadmap that spells out what remains to be proven, especially clinician-annotated localization and a full re-run of the five-method XAI battery on the corrected split (Section~\ref{sec:future-research}). Sixth, open reproducibility. All XAI code ships with this submission, and trained checkpoints and generated figures will go to a permanent public archive once the paper is accepted.

Here is how the rest of the paper unfolds. Section~2 covers background and related work, and Section~3 lays out the materials and methodology. Section~4 presents classification and XAI results. Section~5 discusses what they mean, their clinical relevance, reliability, how they relate to prior work, limitations, and future directions. Section~6 wraps things up.
